\documentclass[11pt,a4paper]{article}

\usepackage[margin=25mm]{geometry}
\usepackage{amsmath,amssymb,bm}
\usepackage{booktabs}
\usepackage{microtype}
\usepackage[hidelinks]{hyperref}
\usepackage{xcolor}
\usepackage{enumitem}

\definecolor{accent}{RGB}{52,73,94}
\hypersetup{colorlinks=true,linkcolor=accent,urlcolor=accent,citecolor=accent}

\newcommand{\R}{\mathbb{R}}
\newcommand{\ind}{\mathbf{1}}
\newcommand{\clamp}{\operatorname{clamp}}
\newcommand{\erf}{\operatorname{erf}}
\newcommand{\Area}{\operatorname{Area}}

\title{Analytically Integrated Additive Saturating Speckle Models\\
\large Disk and Gaussian-Blob Fields over Rectangular Pixels}
\author{}
\date{}

\begin{document}
\maketitle

\begin{abstract}
This note defines two synthetic speckle-pattern models designed for analytically
integrated image generation: an additive disk model and an additive Gaussian-blob
model. Individual blob contributions are integrated over each rectangular pixel,
then summed. Intensity scaling, saturation, and bit-depth quantisation are applied
only after this integration. The construction is therefore not a Boolean union of
disk supports; it is an additive field with a pixel-level saturation nonlinearity.
Its purpose is to provide controlled, deterministic image gradients for digital
image correlation (DIC) while avoiding numerical pixel quadrature.
\end{abstract}

\section{Pixel-box image formation}

Let a rectangular pixel be
\begin{equation}
P=[x_0,x_1]\times[y_0,y_1],
\qquad
|P|=(x_1-x_0)(y_1-y_0).
\end{equation}
A square pixel of pitch $h$ is the special case $x_1-x_0=y_1-y_0=h$.
For a continuous raw texture field $s(\bm{x})$, the exact pixel-box average is
\begin{equation}
\bar{s}_P
=
\frac{1}{|P|}\int_P s(\bm{x})\,\mathrm{d}\bm{x}.
\label{eq:pixel-average}
\end{equation}

The models below use an additive field
\begin{equation}
s(\bm{x})=\sum_{i=1}^{N} s_i(\bm{x}).
\end{equation}
By linearity,
\begin{equation}
\bar{s}_P
=
\frac{1}{|P|}\sum_{i=1}^{N}\int_P s_i(\bm{x})\,\mathrm{d}\bm{x}.
\label{eq:additive-pixel-average}
\end{equation}
Thus each blob may be integrated independently, including when blob supports
overlap.

\subsection{Saturation, intensity mapping, and quantisation}

Each blob has peak raw coverage one.  After all analytically integrated blob
averages have been added, the pixel coverage is saturated:
\begin{equation}
S_P=\clamp\!\left(\bar{s}_P,0,1\right),
\qquad
\clamp(z,0,1)=\min\!\left(1,\max(0,z)\right).
\label{eq:clamp}
\end{equation}
The recorded normalised intensity uses the existing mean and contrast
parameters $I_0$ and $\gamma$:
\begin{equation}
I_P=I_0+\gamma\left(1-2S_P\right).
\label{eq:affine-map}
\end{equation}
For a $B$-bit camera, the quantised integer grey level is
\begin{equation}
g_P
=
\operatorname{round}\!\left[(2^B-1)I_P\right].
\label{eq:quantisation}
\end{equation}
Equations~\eqref{eq:clamp}--\eqref{eq:quantisation} are deliberately applied
\emph{after} analytic pixel integration. In general,
\begin{equation}
\int_P \clamp\!\left(s(\bm{x}),0,1\right)\,\mathrm{d}\bm{x}
\neq
\clamp\!\left(\int_P s(\bm{x})\,\mathrm{d}\bm{x},0,|P|\right).
\end{equation}
The proposed model uses the right-hand ordering as a defined image-generation
rule; it is not intended to reproduce the exact geometric union of opaque paint
disks.

\section{Additive saturating disk model}

\subsection{Raw disk field}

Disk $i$ has centre $\bm{c}_i=(c_{x,i},c_{y,i})$, radius $R_i$, and signed
amplitude $A_i$. Its contribution is
\begin{equation}
s_i^{\mathrm{disk}}(x,y)
=
A_i\,
\ind\!\left[(x-c_{x,i})^2+(y-c_{y,i})^2\le R_i^2\right].
\label{eq:disk-field}
\end{equation}
Positive and negative amplitudes permit bright-on-dark or dark-on-bright
constructions. The complete field is
\begin{equation}
s^{\mathrm{disk}}(x,y)
=
\sum_{i=1}^{N}s_i^{\mathrm{disk}}(x,y).
\end{equation}
Overlap is additive: two coincident unit-amplitude disks produce a raw value of
two before the final intensity mapping and saturation.

\subsection{Exact rectangular-pixel integral}

For disk $i$,
\begin{equation}
\int_P s_i^{\mathrm{disk}}(x,y)\,\mathrm{d}x\,\mathrm{d}y
=
A_i\,\Area(P\cap D_i),
\label{eq:disk-area-basic}
\end{equation}
where $D_i$ is the disk support. Therefore
\begin{equation}
\boxed{
\bar{s}^{\mathrm{disk}}_P
=
\frac{1}{|P|}
\sum_{i=1}^{N}A_i\,\Area(P\cap D_i)
}
\label{eq:disk-pixel-average}
\end{equation}
is exact, apart from floating-point evaluation of elementary functions.

To evaluate the intersection area, translate the pixel into coordinates centred
on disk $i$:
\begin{equation}
a=x_0-c_{x,i},\quad b=x_1-c_{x,i},\quad
c=y_0-c_{y,i},\quad d=y_1-c_{y,i}.
\end{equation}
Define
\begin{equation}
h(u)=\sqrt{R_i^2-u^2},\qquad |u|\le R_i.
\end{equation}
The exact area is the one-dimensional expression
\begin{equation}
\Area(P\cap D_i)
=
\int_{\max(a,-R_i)}^{\min(b,R_i)}
\left[
\min\!\left(d,h(u)\right)
-
\max\!\left(c,-h(u)\right)
\right]_+\,\mathrm{d}u,
\label{eq:disk-slice-integral}
\end{equation}
where $[z]_+=\max(z,0)$. Although written as an integral,
Equation~\eqref{eq:disk-slice-integral} is piecewise analytic. The required
circle antiderivative is
\begin{equation}
H(u;R)
=
\frac{1}{2}
\left[
 u\sqrt{R^2-u^2}
 +R^2\arcsin\!\left(\frac{u}{R}\right)
\right],
\qquad
\frac{\mathrm{d}H}{\mathrm{d}u}=\sqrt{R^2-u^2}.
\label{eq:circle-antiderivative}
\end{equation}

A robust exact implementation proceeds as follows:
\begin{enumerate}[leftmargin=*,itemsep=2pt]
\item Clip the horizontal interval to $[a,b]\cap[-R_i,R_i]$.
\item Form breakpoints from the clipped endpoints and every valid intersection
      of a horizontal pixel edge with the circle:
      \begin{equation}
      u=\pm\sqrt{R_i^2-c^2}\quad (|c|<R_i),
      \qquad
      u=\pm\sqrt{R_i^2-d^2}\quad (|d|<R_i).
      \end{equation}
\item Sort and deduplicate the breakpoints.
\item On each open subinterval, evaluate one midpoint to determine which branches
      of $\min(d,h)$, $\max(c,-h)$, and $[\cdot]_+$ are active.
\item Integrate that fixed branch exactly. Every branch is a linear combination
      of a constant and $\pm h(u)$, so its definite integral uses only endpoint
      differences and $H(u;R_i)$ from Equation~\eqref{eq:circle-antiderivative}.
\end{enumerate}

Useful early-out cases are
\begin{align}
P\cap D_i=\varnothing
&\quad\Rightarrow\quad \Area(P\cap D_i)=0,\\
P\subseteq D_i
&\quad\Rightarrow\quad \Area(P\cap D_i)=|P|,\\
D_i\subseteq P
&\quad\Rightarrow\quad \Area(P\cap D_i)=\pi R_i^2.
\end{align}

\subsection{Why this is not a Boolean union}

A Boolean disk pattern would use
\begin{equation}
s_{\mathrm{bool}}(\bm{x})
=
A\,\ind\!\left[\bm{x}\in\bigcup_iD_i\right],
\end{equation}
whose pixel integral requires the union area
$\Area(P\cap\bigcup_iD_i)$. The additive model instead uses
\begin{equation}
\sum_i A_i\Area(P\cap D_i),
\end{equation}
so overlaps are intentionally counted more than once before saturation. This
choice eliminates multi-disk union geometry while retaining exact integration of
every individual disk.

\section{Additive saturating Gaussian-blob model}

\subsection{Isotropic Gaussian field}

An isotropic Gaussian blob is
\begin{equation}
s_i^{\mathrm{G}}(x,y)
=
A_i\exp\!\left[
-\frac{(x-c_{x,i})^2+(y-c_{y,i})^2}{2\sigma_i^2}
\right],
\label{eq:gaussian-field}
\end{equation}
where $A_i$ is the peak amplitude and $\sigma_i$ is the standard deviation.
The full field is additive:
\begin{equation}
s^{\mathrm{G}}(x,y)=\sum_{i=1}^{N}s_i^{\mathrm{G}}(x,y).
\end{equation}
Because the exponent is separable in $x$ and $y$, the rectangular-pixel integral
factorises exactly.

Define
\begin{align}
E_{x,i}
&=
\erf\!\left(\frac{x_1-c_{x,i}}{\sqrt{2}\sigma_i}\right)
-
\erf\!\left(\frac{x_0-c_{x,i}}{\sqrt{2}\sigma_i}\right),\\
E_{y,i}
&=
\erf\!\left(\frac{y_1-c_{y,i}}{\sqrt{2}\sigma_i}\right)
-
\erf\!\left(\frac{y_0-c_{y,i}}{\sqrt{2}\sigma_i}\right).
\end{align}
Then
\begin{equation}
\int_P s_i^{\mathrm{G}}(x,y)\,\mathrm{d}x\,\mathrm{d}y
=
A_i\frac{\pi\sigma_i^2}{2}E_{x,i}E_{y,i},
\label{eq:gaussian-pixel-integral}
\end{equation}
and hence
\begin{equation}
\boxed{
\bar{s}^{\mathrm{G}}_P
=
\frac{1}{|P|}
\sum_{i=1}^{N}
A_i\frac{\pi\sigma_i^2}{2}E_{x,i}E_{y,i}
}
\label{eq:gaussian-pixel-average}
\end{equation}
without supersampling or numerical quadrature.

The whole-plane integral of one blob is
\begin{equation}
\int_{\R^2}s_i^{\mathrm{G}}(x,y)\,\mathrm{d}x\,\mathrm{d}y
=2\pi A_i\sigma_i^2.
\end{equation}
If instead blob $i$ is normalised to prescribed integrated weight $w_i$, use
\begin{equation}
s_i^{\mathrm{G}}(x,y)
=
\frac{w_i}{2\pi\sigma_i^2}
\exp\!\left[-\frac{\|\bm{x}-\bm{c}_i\|^2}{2\sigma_i^2}\right].
\end{equation}

\subsection{Axis-aligned anisotropic Gaussian blobs}

For widths $\sigma_{x,i}$ and $\sigma_{y,i}$,
\begin{equation}
s_i^{\mathrm{AG}}(x,y)
=
A_i\exp\!\left[
-\frac{(x-c_{x,i})^2}{2\sigma_{x,i}^2}
-\frac{(y-c_{y,i})^2}{2\sigma_{y,i}^2}
\right].
\end{equation}
The exact pixel integral is
\begin{equation}
\int_P s_i^{\mathrm{AG}}\,\mathrm{d}x\,\mathrm{d}y
=
A_i\frac{\pi\sigma_{x,i}\sigma_{y,i}}{2}
E_{x,i}^{(\sigma_x)}E_{y,i}^{(\sigma_y)},
\end{equation}
where each error-function difference uses the corresponding directional width.
The whole-plane integral is $2\pi A_i\sigma_{x,i}\sigma_{y,i}$.

A rotated or correlated Gaussian remains analytically integrable over the whole
plane, but its axis-aligned rectangular-pixel integral is expressed through a
bivariate normal cumulative distribution function rather than a product of two
ordinary error functions.

\section{Optional compact support for Gaussian blobs}

The untruncated Gaussian has infinite support, although its contribution becomes
exponentially small. A circular cutoff at $R_i=k\sigma_i$ defines
\begin{equation}
s_{i,k}^{\mathrm{G}}(\bm{x})
=
s_i^{\mathrm{G}}(\bm{x})
\ind\!\left[\|\bm{x}-\bm{c}_i\|\le k\sigma_i\right].
\end{equation}
Its whole-plane integral is
\begin{equation}
2\pi A_i\sigma_i^2\left(1-e^{-k^2/2}\right).
\end{equation}
However, a circular cutoff destroys Cartesian separability for pixels intersecting
the cutoff boundary. Pixels wholly inside the support retain
Equation~\eqref{eq:gaussian-pixel-integral}; pixels wholly outside contribute zero;
boundary pixels require a piecewise one-dimensional integral or numerical
quadrature. A square cutoff
$|x-c_{x,i}|\le k\sigma_i$, $|y-c_{y,i}|\le k\sigma_i$ preserves separability by
clipping the pixel limits before evaluating the error-function expression.

\section{Comparison of the two models}

\begin{table}[ht]
\centering
\small
\begin{tabular}{@{}lll@{}}
\toprule
Property & Additive disks & Additive Gaussians \\
\midrule
Primitive field & Piecewise constant & Smooth \\
Individual pixel integral & Circle--rectangle area & Product of error functions \\
Overlap handling & Additive before saturation & Additive before saturation \\
Compact support & Exact & Optional approximation/cutoff \\
Boundary regularity & Discontinuous & Infinitely differentiable \\
Numerical quadrature needed & No & No, without circular cutoff \\
Main cost & Piecewise geometry & Special-function evaluation \\
\bottomrule
\end{tabular}
\end{table}

Both constructions provide deterministic pixel values and permit direct control of
blob centres, sizes, amplitudes, density, mean intensity, contrast, saturation, and
bit depth. The Gaussian model offers the simplest closed-form rectangular integral
and smooth image gradients. The disk model provides compact, sharp-edged features
while remaining analytically integrable through exact disk--rectangle intersection
areas.

\section{Recommended image-generation sequence}

For either model, a clear computational definition is
\begin{enumerate}[leftmargin=*,itemsep=2pt]
\item Generate blob parameters $(\bm{c}_i,R_i,A_i)$ or
      $(\bm{c}_i,\sigma_i,A_i)$.
\item Identify blobs whose support or effective support can influence pixel $P$.
\item Evaluate every individual blob--pixel integral analytically.
\item Sum the contributions and divide by pixel area.
\item Apply affine mean/contrast mapping.
\item Saturate to the camera range.
\item Quantise to the requested bit depth.
\end{enumerate}

This ordering isolates the analytically integrated texture model from the later
nonlinear camera operations. It is therefore well suited to renderer-convergence
and DIC uncertainty studies in which pixel-integration error must be separated from
DIC systematic error.

\end{document}
